\noindent\textbf{参数高效与重要性引导的适配。}
参数高效微调减少适配预训练模型时需要更新的参数。LoRA 通过附着于选定权重矩阵的低秩因子表示更新，QLoRA 则结合量化的冻结主干与可训练适配器~\cite{hu2022lora,dettmers2023qlora}。这些机制处理的是适配开销的不同部分。选择目标子集进一步引入“可训练容量放在哪里”的决策，而不仅是容量大小或存储精度。

重要性感知适配已经表明，有限的参数预算不必均匀分配。AdaLoRA 通过对 SVD 风格的更新参数化进行重要性引导的剪枝，在权重矩阵之间分配非均匀秩预算~\cite{fm_adalora2023}。FISH 使用基于 Fisher 的评分选择稀疏可训练参数掩码~\cite{fm_fish2021}。FedMATRIX-LoRA 在另一分配层次上延续重要性引导的选择：为各联邦客户端选择固定秩目标子集。重要性感知秩分配控制权重更新内部的容量，本文的形式则在异构客户端和轮次级约束下协调可独立选择的目标投影。

\noindent\textbf{异构联邦 LoRA 分配。}
HetLoRA 通过本地秩自剪枝和稀疏加权聚合支持不同的客户端秩，FAH-QLoRA 则联合考虑量化与 LoRA 秩~\cite{cho2024hetlora,gao2025fah}。这些方法根据客户端资源调整容量或数值精度。Fed-HeLLo 指出另一个内存问题：在其研究的激活占主导设置中，降低秩主要影响优化器相关状态，因此部分层训练具有动机~\cite{zhang2025fedhello}。LEGEND 联合配置微调深度和秩分布，明确将可训练位置与资源使用联系起来~\cite{liu2025legend}。因此，目标选择与秩分配是互补而非互斥的维度。

Fed-HeLLo 通过服务器代理评分和几何分配策略考虑层重要性的差异~\cite{zhang2025fedhello}。其评分引导策略将共享评估转化为分配概率，因此客户端可以获得不同的层掩码。FedMATRIX-LoRA 则允许分别选择块内投影，并在客户端数据上评估评分准则。图~\ref{fig:fm-motivation} 的受控对比将矩阵级和层级选择同时应用于两种评分来源，从而区分粒度问题与评分来源问题。它比较的是共同流程中的分配约束，并非对 Fed-HeLLo 的完整复现。

Fed-pilot 是密切相关的重要性引导分配方法。其按块索引的形式结合本地及汇总的梯度信息增益评分、内存约束分配和补偿聚合~\cite{fm_fedpilot2025}。FedMATRIX-LoRA 在块内开放单独的目标投影指示变量，并在内存上限、时延约束和共享轮次预算下跨客户端协调配置。相关差异在于分配决策及其耦合约束，而不只是本地评分或背包求解。

\noindent\textbf{聚合与资源协调。}
FedEx-LoRA 处理聚合低秩因子时的不匹配，LoRA-FAIR 则改进聚合和初始化~\cite{singhal2025fedex,bian2025lorafair}。这些研究关注如何结合本地更新。Oort、FedBalancer 和 PyramidFL 处理异构条件下的参与者选择或训练工作协调~\cite{lai2021oort,shin2022fedbalancer,li2022pyramidfl}。FedMATRIX-LoRA 将可训练目标子集纳入分配给参与客户端的配置。配置生成将本地更新优先级连接到资源可行性，服务器协调决定哪些配置可一起执行。按矩阵收敛性分析刻画由这一过程产生的选择性更新模式。
